% Ad Quintum Fratrem 2.7 (ad-quintum-fratrem-02-07)
% Source: https://raw.githubusercontent.com/PerseusDL/canonical-latinLit/master/data/phi0474/phi058/phi0474.phi058.perseus-lat1.xml
% Fetched by scripts/fetch_latin.py

% Dateline (Perseus): Scr. Romae mense Februado a. 699 (55) .

\ciceroLetterOpener{MARCVS QVINTO FRATRI SALVTEM.}

\ciceroSection{1} placiturum tibi esse librum meum suspicabar ; tam valde placuisse quam scribis valde gaudeo. quod me admones de nostra Vrania suadesque ut meminerim Iovis orationem quae est in extremo illo libro, ego vero memini et illa omnia mihi magis scripsi quam ceteris.

\ciceroSection{2} sed tamen postridie quam tu es profectus multa nocte cum Vibullio veni ad Pompeium ; cumque ego egissem de istis operibus atque inscriptionibus, per mihi benigne respondit magnam spem attulit ; cum Crasso se dixit loqui velle mihique ut idem facerem suasit. Crassum consulem ex senatu domum reduxi. suscepit rem dixitque esse quod Clodius hoc tempore cuperet per se et per Pompeium consequi ; putare se, si ego eum non impedirem, posse me adipisci sine contentione quod vellem. totum ei negotium permisi meque in eius potestate dixi fore. interfuit huic sermoni P. Crassus adulescens nostri, ut scis, studiosissimus. illud autem quod cupit Clodius est legatio aliqua (si minus per senatum, per populum) libera aut Byzantium aut ad Brogitarum aut utrumque plena res nummorum. quod ego non nimium laboro, etiam si minus adsequor quod volo. Pompeius tamen cum Crasso locutus est. videntur negotium suscepisse. si perficiunt, optime ; si minus, ad nostrum Iovem revertamur.

\ciceroSection{3} A. d. iii Idus Febr. senatus consultum est factum de a ambitu in Afrani sententiam, quam ego dixeram cum tu adesses ; sed magno cum gemitu senatus consules non sunt persecuti eorum sententias qui, Afranio cum essent adsensi, addiderunt ut praetores ita crearentur ut dies sexaginta privati essent. eo die Catonem plane repudiarunt. quid multa? tenent omnia idque ita omnis intellegere volunt.
